\begin{document}
\selectlanguage{brazil}

% Capa (elemento externo — não entra na contagem de páginas)
\imprimircapa

% Folha de rosto (inicia a contagem sem exibir número)
\imprimirfolhaderosto*

% Ficha catalográfica (substituir pelo PDF emitido pela BU/UFSC antes da entrega final)
\begin{fichacatalografica}
      \vspace*{\fill}
      \begin{center}
            \fbox{%
                  \begin{minipage}{12.5cm}
                        \vspace{0.6cm}
                        \centering
                        \textbf{Ficha Catalográfica}\\[0.5cm]
                        Espaço reservado para a ficha catalográfica gerada pelo sistema da
                        Biblioteca Universitária da UFSC (\url{https://portal.bu.ufsc.br/}).\\
                        Quando disponível, substitua este bloco por \texttt{\textbackslash includepdf\{ficha\_catalografica.pdf\}}.
                        \vspace{0.6cm}
                  \end{minipage}%
            }
      \end{center}
\end{fichacatalografica}
\clearpage

% Folha de aprovação (preencher data e membros da banca antes da versão definitiva)
\begin{folhadeaprovacao}
      \begin{center}
            {\ABNTEXchapterfont\large\imprimirautor}

            \vspace*{\fill}
            {\ABNTEXchapterfont\bfseries\Large\imprimirtitulo}
            \vspace*{\fill}
      \end{center}

      \hspace{.45\textwidth}
      \begin{minipage}{.5\textwidth}
            Este Trabalho de Conclusão de Curso foi avaliado e aprovado como requisito
            parcial para a obtenção do título de Bacharel em Engenharia Mecânica pela
            Universidade Federal de Santa Catarina.
      \end{minipage}

      \vspace{\onelineskip}
      \begin{center}
            \imprimirlocal, [dia] de [mês] de \imprimirdata.
      \end{center}

      \vspace*{\fill}
      \assinatura{\textbf{\imprimirorientador} \\ Orientador \\ Universidade Federal de Santa Catarina}
      \assinatura{\textbf{Prof.~[Membro da Banca 1]} \\ Avaliador \\ Universidade Federal de Santa Catarina}
      \assinatura{\textbf{Prof.~[Membro da Banca 2]} \\ Avaliador \\ Universidade Federal de Santa Catarina}
      \vspace*{\fill}
\end{folhadeaprovacao}

\begin{resumo}\label{resumo}
      Algoritmos meméticos que combinam Algoritmos Genéticos com busca local \(2\)-opt produzem soluções competitivas para o Problema do Caixeiro Viajante (TSP), mas impõem elevado custo computacional conforme o número de cidades cresce.
      Este trabalho investiga o impacto de diferentes estratégias de paralelização em Unidade de Processamento Gráfico (GPU) sobre o desempenho desse método, avaliando como a divisão de tarefas entre CPU e GPU e o tráfego de dados pelo barramento PCIe afetam o tempo de execução e a qualidade das soluções.
      Para isso, foram implementadas e comparadas quatro variantes sobre o mesmo algoritmo base e os mesmos critérios de parada: execução integral em CPU, arquitetura híbrida com envio individual de rotas para a GPU, arquitetura híbrida com envio conjunto da população e execução totalmente residente em GPU.
      Os experimentos foram conduzidos em \totalinstancias\ instâncias da biblioteca TSPLIB, de \menorinstancia\ a \maiorinstancia\ cidades, com \totalrepeticoes\ repetições independentes e avaliação estatística não paramétrica.
      Os resultados mostram que todas as quatro variantes escolhidas para a comparação atingem soluções próximas do ótimo conhecido sem diferença estatisticamente significativa entre si (\emph{gaps} médios entre \(0{,}90\%\) e \(0{,}91\%\); teste de Friedman, \(p = 0{,}4543\)), enquanto a variante totalmente residente em GPU obtém os menores tempos de execução em todas as instâncias avaliadas, alcançando ganho médio geométrico de \(71{,}03\times\) frente à CPU, contra \(23{,}86\times\) da híbrida otimizada e \(4{,}57\times\) da híbrida ingênua.
      Conclui-se que reduzir transferências entre CPU e GPU e estruturar adequadamente os \emph{kernels} na memória do dispositivo é determinante para transformar o paralelismo da GPU em ganho efetivo de velocidade sem perda de qualidade.

      \vspace{\onelineskip}
      \noindent\textbf{Palavras-chave:} Problema do Caixeiro Viajante; Algoritmos Genéticos; Computação em GPU; CUDA; Otimização Combinatória.
\end{resumo}

\begin{resumo}[Abstract]\label{abstract}
      Memetic algorithms combining Genetic Algorithms with \(2\)-opt local search produce competitive solutions for the Traveling Salesman Problem (TSP), but incur high computational cost as the number of cities grows.
      This work investigates the impact of different Graphics Processing Unit (GPU) parallelization strategies on the performance of this method, examining how workload partitioning between CPU and GPU and data traffic across the PCIe bus affect execution time and solution quality.
      Four variants were implemented and compared using the same base algorithm and stopping criteria: a CPU-only reference, a hybrid variant transferring routes individually to the GPU, a hybrid variant transferring the entire population jointly, and a fully GPU-resident variant.
      Experiments were carried out on \totalinstancias\ TSPLIB instances ranging from \menorinstancia\ to \maiorinstancia\ cities, with \totalrepeticoes\ independent runs and non-parametric statistical evaluation.
      The results show that all four chosen variants achieve solutions close to the known optimum with no statistically significant difference among them (mean gaps between \(0.90\%\) and \(0.91\%\); Friedman test, \(p = 0.4543\)), whereas the fully GPU-resident variant delivers the lowest runtimes across all evaluated instances, achieving a geometric mean speedup of \(71.03\times\) over the CPU baseline, compared to \(23.86\times\) for the optimized hybrid and \(4.57\times\) for the naive hybrid.
      These findings confirm that minimizing CPU--GPU data transfers and designing kernels suited to the device memory hierarchy are decisive for turning GPU parallelism into practical speedup without sacrificing solution quality.

      \vspace{\onelineskip}
      \noindent\textbf{Keywords:} Traveling Salesman Problem; Genetic Algorithms; GPU Computing; CUDA; Combinatorial Optimization.
\end{resumo}

\pdfbookmark[0]{\listfigurename}{lof}
\listoffigures*
\clearpage

\pdfbookmark[0]{\listtablename}{lot}
\listoftables*
\clearpage

\pretextualchapter{\listadesiglasname}\label{lista-de-siglas}

\begin{longtable}[]{@{}
            >{\raggedright\arraybackslash}p{(\linewidth - 2\tabcolsep) * \real{0.20}}
            >{\raggedright\arraybackslash}p{(\linewidth - 2\tabcolsep) * \real{0.80}}@{}}
      \toprule\noalign{}
      Sigla / Abreviatura & Significado                                                                                             \\
      \midrule\noalign{}
      \endhead
      \bottomrule\noalign{}
      \endlastfoot
      ACO                 & Otimização por Colônia de Formigas (\emph{Ant Colony Optimization})                                     \\
      \sepitem
      AG                  & Algoritmo Genético (em inglês, \emph{Genetic Algorithm} --- GA)                                         \\
      \sepitem
      ALU                 & Unidade Lógica e Aritmética (\emph{Arithmetic Logic Unit})                                              \\
      \sepitem
      ANOVA               & Análise de Variância (\emph{Analysis of Variance})                                                      \\
      \sepitem
      API                 & Interface de Programação de Aplicações (\emph{Application Programming Interface})                       \\
      \sepitem
      ATSP                & Problema do Caixeiro Viajante Assimétrico (\emph{Asymmetric Traveling Salesman Problem})                \\
      \sepitem
      AVX                 & Extensões Vetoriais Avançadas (\emph{Advanced Vector Extensions})                                       \\
      \sepitem
      CD                  & Diferença Crítica (\emph{Critical Difference})                                                          \\
      \sepitem
      CNC                 & Controle Numérico Computadorizado (\emph{Computer Numerical Control})                                   \\
      \sepitem
      CPU                 & Unidade Central de Processamento (\emph{Central Processing Unit})                                       \\
      \sepitem
      CUDA                & Arquitetura Unificada de Dispositivos de Computação (\emph{Compute Unified Device Architecture})        \\
      \sepitem
      CVRP                & Problema de Roteamento de Veículos Capacitado (\emph{Capacitated Vehicle Routing Problem})              \\
      \sepitem
      D2H                 & Transferência do dispositivo para o hospedeiro (\emph{Device-to-Host}, GPU \(\to\) CPU)                 \\
      \sepitem
      DRAM                & Memória Dinâmica de Acesso Aleatório (\emph{Dynamic Random-Access Memory})                              \\
      \sepitem
      EAX                 & Cruzamento por Montagem de Arestas (\emph{Edge Assembly Crossover})                                     \\
      \sepitem
      FLS                 & Busca Local Rápida (\emph{Fast Local Search})                                                           \\
      \sepitem
      GCC                 & Coleção de Compiladores GNU (\emph{GNU Compiler Collection})                                            \\
      \sepitem
      GLS                 & Busca Local Guiada (\emph{Guided Local Search})                                                         \\
      \sepitem
      GPGPU               & Computação de Propósito Geral em GPUs (\emph{General-Purpose Computing on GPUs})                        \\
      \sepitem
      GPU                 & Unidade de Processamento Gráfico (\emph{Graphics Processing Unit})                                      \\
      \sepitem
      H2D                 & Transferência do hospedeiro para o dispositivo (\emph{Host-to-Device}, CPU \(\to\) GPU)                 \\
      \sepitem
      HPC                 & Computação de Alto Desempenho (\emph{High Performance Computing})                                       \\
      \sepitem
      IIQ                 & Intervalo Interquartil (\(Q_{25}\text{--}Q_{75}\))                                                      \\
      \sepitem
      JSON                & Notação de Objetos JavaScript (\emph{JavaScript Object Notation})                                       \\
      \sepitem
      LKH                 & Heurística de Lin--Kernighan--Helsgaun (\emph{Lin--Kernighan--Helsgaun})                                \\
      \sepitem
      MDVRP               & Problema de Roteamento de Veículos com Múltiplos Depósitos (\emph{Multi-Depot Vehicle Routing Problem}) \\
      \sepitem
      MIMD                & Múltiplas Instruções, Múltiplos Dados (\emph{Multiple Instruction, Multiple Data})                      \\
      \sepitem
      MST                 & Árvore Geradora Mínima (\emph{Minimum Spanning Tree})                                                   \\
      \sepitem
      NN                  & Heurística do Vizinho Mais Próximo (\emph{Nearest Neighbor})                                            \\
      \sepitem
      OX                  & Cruzamento por Ordem (\emph{Order Crossover})                                                           \\
      \sepitem
      PCIe                & Interconexão de Componentes Periféricos Expressa (\emph{Peripheral Component Interconnect Express})     \\
      \sepitem
      SIMD                & Instrução Única, Múltiplos Dados (\emph{Single Instruction, Multiple Data})                             \\
      \sepitem
      SIMT                & Instrução Única, Múltiplas \emph{Threads} (\emph{Single Instruction, Multiple Threads})                 \\
      \sepitem
      SM                  & Multiprocessador de Fluxo (\emph{Streaming Multiprocessor})                                             \\
      \sepitem
      SSE                 & Extensões SIMD de Fluxo (\emph{Streaming SIMD Extensions})                                              \\
      \sepitem
      TSP                 & Problema do Caixeiro Viajante (\emph{Traveling Salesman Problem})                                       \\
      \sepitem
      TSPLIB              & Biblioteca de Instâncias de Referência para o TSP (\emph{Traveling Salesman Problem Library})           \\
      \sepitem
      VRAM                & Memória de Vídeo de Acesso Aleatório da GPU (\emph{Video Random-Access Memory})                         \\
\end{longtable}

\clearpage

\pretextualchapter{Lista de Símbolos}\label{lista-de-simbolos}

\begin{longtable}[]{@{}
            >{\raggedright\arraybackslash}p{(\linewidth - 4\tabcolsep) * \real{0.3333}}
            >{\raggedright\arraybackslash}p{(\linewidth - 4\tabcolsep) * \real{0.3333}}
            >{\raggedright\arraybackslash}p{(\linewidth - 4\tabcolsep) * \real{0.3333}}@{}}
      \toprule\noalign{}
      Símbolo                                             & Descrição                                                                                                                                     & Ver seção \\
      \midrule\noalign{}
      \endhead
      \bottomrule\noalign{}
      \endlastfoot
      \(G = (V, E)\)                                      & Grafo completo não direcionado composto por vértices \(V\) e arestas \(E\)                                                                    &
      \ref{sec-21-otimizacao-tsp}                                                                                                                                                                                     \\
      \sepitem
      \(V\)                                               & Conjunto de vértices (cidades) do grafo                                                                                                       &
      \ref{sec-21-otimizacao-tsp}                                                                                                                                                                                     \\
      \sepitem
      \(E\)                                               & Conjunto de arestas entre pares de cidades                                                                                                    &
      \ref{sec-21-otimizacao-tsp}                                                                                                                                                                                     \\
      \sepitem
      \(c_{ij}\)                                          & Custo (distância) entre as cidades \(i\) e \(j\)                                                                                              &
      \ref{sec-21-otimizacao-tsp} e \ref{sec-221-2opt}                                                                                                                                                                \\
      \sepitem
      \(n_{coords}\) (ou \(n\))                           & Número de cidades (coordenadas) da instância                                                                                                  &
      \ref{sec-21-otimizacao-tsp}, \ref{sec-221-2opt}, \ref{sec-311-hardware-software} e \ref{sec-33-selecao-instancias}                                                                                              \\
      \sepitem
      \(n_{\text{inst}}\)                                 & Número de instâncias da TSPLIB avaliadas (\(n_{\text{inst}} = \totalinstancias\))                                                             &
      \ref{sec-33-selecao-instancias} e \ref{sec-34-protocolo-estatistico}                                                                                                                                            \\
      \sepitem
      \(m_{\text{alg}}\)                                  & Número de variantes algorítmicas comparadas (\(m_{\text{alg}} = \totalvariantes\))                                                            &
      \ref{sec-34-protocolo-estatistico} e \ref{sec-43-analise-agregada}                                                                                                                                              \\
      \sepitem
      \(n_{pop}\)                                         & Tamanho da população do Algoritmo Genético (\(n_{pop} = 256\))                                                                                &
      \ref{sec-23-algoritmos-geneticos} e \ref{sec-321-ga-base}                                                                                                                                                       \\
      \sepitem
      \(\mu, \lambda\)                                    & Tamanhos da população de pais (\(\mu\)) e de descendentes (\(\lambda\)) na seleção elitista \((\mu + \lambda)\) (\(\mu = \lambda = n_{pop}\)) &
      \ref{sec-231-operadores-substituicao}, \ref{sec-251-ga-2opt}, \ref{sec-321-ga-base} e \ref{apendice-pseudocodigos}                                                                                              \\
      \sepitem
      \(n_{reps}\)                                        & Número de repetições independentes por experimento (\(n_{reps} = \totalrepeticoes\))                                                          &
      \ref{sec-14-objetivos-especificos}, \ref{sec-311-hardware-software}, \ref{sec-33-selecao-instancias} e \ref{sec-34-protocolo-estatistico}                                                                       \\
      \sepitem
      \(p_c\)                                             & Taxa de cruzamento (\(p_c = 1{,}0\))                                                                                                          &
      \ref{sec-231-operadores-substituicao}, \ref{sec-251-ga-2opt}, \ref{sec-321-ga-base} e \ref{apendice-pseudocodigos}                                                                                              \\
      \sepitem
      \(p_{\text{mut}}\)                                  & Taxa de mutação por rota (\(p_{\text{mut}} = {,}02\))                                                                                         &
      \ref{sec-231-operadores-substituicao}, \ref{sec-251-ga-2opt}, \ref{sec-321-ga-base} e \ref{apendice-pseudocodigos}                                                                                              \\
      \sepitem
      \(\varepsilon_{\text{alvo}}\)                       & Tolerância alvo de \emph{gap} para critério de parada antecipada (\(\varepsilon_{\text{alvo}} = 1\%\))                                        &
      \ref{sec-321-ga-base} e \ref{apendice-pseudocodigos}                                                                                                                                                            \\
      \sepitem
      \(G_{\text{estag}}\)                                & Limite de gerações consecutivas sem melhoria (\(G_{\text{estag}} = \lceil 2\sqrt{n_{coords}} \rceil\))                                        &
      \ref{sec-321-ga-base} e \ref{apendice-pseudocodigos}                                                                                                                                                            \\
      \sepitem
      \(G_{\max}\)                                        & Teto máximo de gerações do ciclo evolutivo (\(G_{\max} = \lceil 2\,n_{coords}\sqrt{n_{coords}} \rceil\))                                      &
      \ref{sec-321-ga-base} e \ref{apendice-pseudocodigos}                                                                                                                                                            \\
      \sepitem
      \(\tau_{\text{trnmt}}\)                             & Tamanho do torneio na seleção de progenitores (\(\tau_{\text{trnmt}} = 5\))                                                                   &
      \ref{sec-231-operadores-substituicao}, \ref{sec-251-ga-2opt}, \ref{sec-321-ga-base} e \ref{apendice-pseudocodigos}                                                                                              \\
      \sepitem
      \(\kappa\)                                          & Limite máximo de varreduras (iterações) da busca local \(2\)-opt (\(\kappa = \maxpassadas\))                                                  &
      \ref{sec-221-2opt}, \ref{sec-24-metaheuristicas-hibridas}, \ref{sec-251-ga-2opt} e \ref{sec-321-ga-base}                                                                                                        \\
      \sepitem
      \(k_{\text{viz}}\)                                  & Número de vizinhos na lista estática de candidatos da busca local (\(k_{\text{viz}} = \kviz\))                                                &
      \ref{sec-221-2opt}, \ref{sec-251-ga-2opt}, \ref{sec-321-ga-base} e \ref{sec-5311-semantica-operadores}                                                                                                          \\
      \sepitem
      \(N_k\)                                             & Matriz estática dos \(k_{\text{viz}}\) vizinhos mais próximos por cidade (\(n_{coords} \times k_{\text{viz}}\))                               &
      \ref{sec-322-variantes} e \ref{sec-44-decomposicao-fases}                                                                                                                                                       \\
      \sepitem
      \(\mathcal{L}(x, \kappa)\)                          & Busca local truncada aplicada à solução \(x\) com limite de \(\kappa\) varreduras                                                             &
      \ref{sec-251-ga-2opt}                                                                                                                                                                                           \\
      \sepitem
      \(C\)                                               & Circuito hamiltoniano (vetor de rota candidata)                                                                                               &
      \ref{sec-221-2opt}, \ref{sec-322-variantes} e \ref{apendice-pseudocodigos}                                                                                                                                      \\
      \sepitem
      \(B, B^*\)                                          & Custo de uma solução heurística \(B\) e custo ótimo \(B^*\)                                                                                   &
      \ref{sec-22-heuristicas-tsp}, \ref{sec-321-ga-base}, \ref{sec-33-selecao-instancias} e \ref{sec-34-protocolo-estatistico}                                                                                       \\
      \sepitem
      \(D\)                                               & Matriz simétrica de distâncias inteiras (\(n_{coords} \times n_{coords}\))                                                                    &
      \ref{sec-311-hardware-software} e \ref{sec-312-ambiente-software}                                                                                                                                               \\
      \sepitem
      \(\Delta C\)                                        & Variação de custo em um movimento local \(2\)-opt                                                                                             &
      \ref{sec-221-2opt}, \ref{sec-312-ambiente-software} e \ref{apendice-pseudocodigos}                                                                                                                              \\
      \sepitem
      \(f(x)\)                                            & Função de aptidão (\emph{fitness}) de uma solução \(x\)                                                                                       &
      \ref{sec-23-algoritmos-geneticos}, \ref{sec-231-operadores-substituicao} e \ref{sec-251-ga-2opt}                                                                                                                \\
      \sepitem
      \(P(t), P_{\mathrm{desc}}\)                         & População na geração \(t\) e conjunto de descendentes gerados (matrizes \(n_{pop} \times n_{coords}\))                                        &
      \ref{sec-23-algoritmos-geneticos}, \ref{sec-231-operadores-substituicao}, \ref{sec-251-ga-2opt}, \ref{sec-322-variantes} e \ref{apendice-pseudocodigos}                                                         \\
      \sepitem
      \(t\)                                               & Índice da geração evolutiva                                                                                                                   &
      \ref{sec-23-algoritmos-geneticos} e \ref{sec-251-ga-2opt}                                                                                                                                                       \\
      \sepitem
      \(T\) (\(T_{total}, T_{cpu}, T_{gpu}, T_{transf}\)) & Tempo de execução (total, em CPU, em GPU e de transferência PCIe)                                                                             &
      \ref{sec-311-hardware-software}, \ref{sec-34-protocolo-estatistico} e \ref{sec-44-decomposicao-fases}                                                                                                           \\
      \sepitem
      \(Mem_{total}\)                                     & Consumo total estimado de VRAM na variante residente                                                                                          &
      \ref{sec-311-hardware-software}                                                                                                                                                                                 \\
      \sepitem
      \(\mathit{H2D}\)                                    & Fluxo de transferência \emph{Host-to-Device} (CPU \(\to\) GPU)                                                                                &
      \ref{sec-14-objetivos-especificos}, \ref{sec-322-variantes} e \ref{sec-44-decomposicao-fases}                                                                                                                   \\
      \sepitem
      \(\mathit{D2H}\)                                    & Fluxo de transferência \emph{Device-to-Host} (GPU \(\to\) CPU)                                                                                &
      \ref{sec-14-objetivos-especificos}, \ref{sec-322-variantes} e \ref{sec-44-decomposicao-fases}                                                                                                                   \\
      \sepitem
      \(H_0\)                                             & Hipótese nula de equivalência de desempenho entre os algoritmos comparados                                                                    &
      \ref{sec-27-comparacao-estatistica}, \ref{sec-272-comparacoes-multiplas} e \ref{sec-34-protocolo-estatistico}                                                                                                   \\
      \sepitem
      \(\alpha\)                                          & Nível nominal de significância estatística (\(\alpha = 0{,}05\))                                                                              &
      \ref{sec-27-comparacao-estatistica}, \ref{sec-34-protocolo-estatistico} e \ref{sec-43-analise-agregada}                                                                                                         \\
      \sepitem
      \(p\) (\(p_F, p_{\text{Holm}}\))                    & Valor-\(p\) de teste estatístico, da estatística \(F_F\) e ajustado por Holm--Bonferroni                                                      &
      \ref{sec-27-comparacao-estatistica}, \ref{sec-34-protocolo-estatistico} e \ref{sec-43-analise-agregada}                                                                                                         \\
      \sepitem
      \(\chi^2_F\)                                        & Estatística qui-quadrado do teste de Friedman                                                                                                 &
      \ref{sec-272-comparacoes-multiplas}, \ref{sec-34-protocolo-estatistico} e \ref{sec-43-analise-agregada}                                                                                                         \\
      \sepitem
      \(F_F\)                                             & Estatística \(F\) corrigida de Iman--Davenport                                                                                                &
      \ref{sec-272-comparacoes-multiplas}, \ref{sec-34-protocolo-estatistico} e \ref{sec-43-analise-agregada}                                                                                                         \\
      \sepitem
      \(CD\)                                              & Diferença crítica (\emph{Critical Difference}) do pós-teste de Nemenyi                                                                        &
      \ref{sec-272-comparacoes-multiplas}, \ref{sec-34-protocolo-estatistico} e \ref{sec-43-analise-agregada}                                                                                                         \\
      \sepitem
      \(R_j, \Delta R\)                                   & Posto médio da \(j\)-ésima variante no teste de Friedman e diferença entre postos médios no pós-teste de Nemenyi                              &
      \ref{sec-272-comparacoes-multiplas}, \ref{sec-34-protocolo-estatistico} e \ref{sec-43-analise-agregada}                                                                                                         \\
      \sepitem
      \(q_\alpha\)                                        & Valor crítico da amplitude studentizada dividida por \(\sqrt{2}\)                                                                             &
      \ref{sec-272-comparacoes-multiplas} e \ref{sec-34-protocolo-estatistico}                                                                                                                                        \\
      \sepitem
      \(W\)                                               & Estatística do teste de postos sinalizados de Wilcoxon                                                                                        &
      \ref{sec-271-comparacoes-pareadas}, \ref{sec-34-protocolo-estatistico} e \ref{sec-43-analise-agregada}                                                                                                          \\
      \sepitem
      \(\sigma\)                                          & Desvio padrão amostral dos resultados                                                                                                         &
      \ref{sec-42-analise-por-instancia}, \ref{sec-43-analise-agregada} e \ref{apendice-resultados}                                                                                                                   \\
      \sepitem
      \(Q_{25}, Q_{75}\)                                  & Primeiro e terceiro quartis (faixa interquartil \(Q_{25}\text{--}Q_{75}\))                                                                    &
      \ref{sec-43-analise-agregada}                                                                                                                                                                                   \\
\end{longtable}
% Zera o contador de tabelas incrementado pelos longtables pré-textuais (Lista de Abreviaturas e Siglas e Lista de Símbolos)
\setcounter{table}{0}

\clearpage

\pdfbookmark[0]{\contentsname}{toc}
\tableofcontents*
